\begin{lemma}[Last crossing-window coefficients from a block-diagonal chain vector]
    \label{lem:block_diagonal_last_crossing_coefficients_from_chain_ground_space}
    \lean{MPSTensor.blockDiagonal_boundary_last_coefficients_of_blockDiagonal_chainGroundSpace}
    \leanok
    \uses{lem:block_diagonal_boundary_local_sum_mem,
        lem:block_diagonal_last_crossing_coefficients_from_sum_membership}
    Let $B=\bigoplus_j\mu_jA^j$, and assume $\mu_j\ne0$ for every block
    $j$. Suppose
    \[
        \psi=\Gamma_{M+1}^{B}\!\left(\bigoplus_jX_j\right),
        \qquad
        \psi\in\mathcal G_{M+1,m+2}(B).
    \]
    Assume also that the length-$m$ simultaneous word tuples span the product
    algebra. Then, for every outside configuration $\tau$, with
    $\rho=\tau_{m+1}\cdots\tau_{M-1}$, there are matrices $E_j$ such that
    \[
        A^j_b\bigl(A^j_\rho A^j_a(\mu_j^{M+1}X_j)\bigr)
        =
        A^j_bE_jA^j_a
    \]
    for every block $j$ and all boundary letters $a,b$.
\end{lemma}

\begin{proof}
    \leanok
    \uses{lem:block_diagonal_boundary_local_sum_mem,
        lem:block_diagonal_last_crossing_coefficients_from_sum_membership}
    The local constraint of $\psi$ at the cyclic interval beginning at $M$
    gives
    \[
        \sum_j
        R_{M,\tau}\!\left(\Gamma_{M+1}^{A_j}(\mu_j^{M+1}X_j)\right)
        \in \bigvee_j G_{m+2}(A^j)
    \]
    by Lemma~\ref{lem:block_diagonal_boundary_local_sum_mem}. Applying
    Lemma~\ref{lem:block_diagonal_last_crossing_coefficients_from_sum_membership}
    gives the displayed equations.
\end{proof}

\begin{lemma}[Last crossing-window component membership from a chain vector]
    \label{lem:block_diagonal_last_crossing_component_membership_from_chain_ground_space}
    \lean{MPSTensor.blockDiagonal_boundary_last_component_mem_groundSpace_of_blockDiagonal_chainGroundSpace}
    \leanok
    \uses{lem:block_diagonal_boundary_local_sum_mem,
        lem:block_diagonal_last_crossing_component_membership_from_sum_membership}
    Let $B=\bigoplus_j\mu_jA^j$, and assume $\mu_j\ne0$ for every block
    $j$. Suppose
    \[
        \psi=\Gamma_{M+1}^{B}\!\left(\bigoplus_jX_j\right),
        \qquad
        \psi\in\mathcal G_{M+1,m+2}(B).
    \]
    If the length-$m$ simultaneous word tuples span the product algebra, then
    for every outside configuration $\tau$ and every block $j$,
    \[
        R_{M,\tau}\!\left(\Gamma_{M+1}^{A_j}(\mu_j^{M+1}X_j)\right)
        \in G_{m+2}(A^j).
    \]
\end{lemma}

\begin{proof}
    \leanok
    \uses{lem:block_diagonal_boundary_local_sum_mem,
        lem:block_diagonal_last_crossing_component_membership_from_sum_membership}
    Lemma~\ref{lem:block_diagonal_boundary_local_sum_mem} applied to the
    window beginning at the last site gives
    \[
        \sum_j
        R_{M,\tau}\!\left(\Gamma_{M+1}^{A_j}(\mu_j^{M+1}X_j)\right)
        \in \bigvee_j G_{m+2}(A^j).
    \]
    Lemma~\ref{lem:block_diagonal_last_crossing_component_membership_from_sum_membership}
    then applies to each block component.
\end{proof}

\begin{lemma}[Boundary-crossing intervals from a right boundary-matrix identity]
    \label{lem:block_diagonal_boundary_crossing_right_boundary_matrix}
    \lean{MPSTensor.blockDiagonal_boundary_cyclicRestrict_component_mem_groundSpace_of_boundary_identity}
    \leanok
    \uses{lem:block_diagonal_boundary_crossing_matrix_identity}
    Assume $0<N$, $L\le N$, and let the cyclic interval beginning at $i$
    cross the boundary cut, so $N<i+L$. Put $a=i+L-N$. Suppose that there is
    a matrix $Y$ such that for every word $\beta$ of length $a$,
    \[
        \mu_j^N X_j A^j_\beta=A^j_\beta Y .
    \]
    Then
    \[
        R_{i,\tau}\!\left(\Gamma_N^{A_j}(\mu_j^NX_j)\right)
        \in G_L(A_j).
    \]
\end{lemma}

\begin{proof}
    \leanok
    \uses{lem:block_diagonal_boundary_crossing_matrix_identity}
    Let
    \[
        E:=Y A^j_{\tau_a}\cdots A^j_{\tau_{i-1}} .
    \]
    Then, for every word $\beta$ on the segment $0,\ldots,a-1$,
    \[
        \mu_j^N X_j A^j_\beta
        A^j_{\tau_a}\cdots A^j_{\tau_{i-1}}
        =
        A^j_\beta E.
    \]
    Lemma~\ref{lem:block_diagonal_boundary_crossing_matrix_identity} gives the
    stated local constraint.
\end{proof}

\begin{lemma}[Boundary-crossing cyclic intervals suffice]
    \label{lem:block_diagonal_boundary_component_crossing_windows_suffice}
    \lean{MPSTensor.blockDiagonal_boundary_component_chainGroundSpace_of_crossing_windows}
    \leanok
    \uses{def:chain_ground_space,
        lem:block_diagonal_boundary_nonwrapping_component}
    Assume $0<N$ and $L\le N$. Fix block-diagonal boundary conditions
    $X_j$. Suppose that, for every block $j$, every cyclic interval
    beginning at $i$ with $N<i+L$, and every configuration $\tau$,
    \[
        R_{i,\tau}\!\left(\Gamma_N^{A_j}(\mu_j^NX_j)\right)
        \in G_L(A_j).
    \]
    Then, for every block $j$,
    \[
        \Gamma_N^{A_j}(\mu_j^NX_j)\in \mathcal G_{N,L}(A_j).
    \]
\end{lemma}

\begin{proof}
    \leanok
    \uses{def:chain_ground_space,
        lem:block_diagonal_boundary_nonwrapping_component}
    By definition, membership in $\mathcal G_{N,L}(A_j)$ is the collection of
    all cyclic length-$L$ local constraints. For an interval beginning at $i$,
    either $i+L\le N$ or $N<i+L$. In the first case use
    Lemma~\ref{lem:block_diagonal_boundary_nonwrapping_component}. In the second
    case use the displayed hypothesis.
\end{proof}

\begin{lemma}[Right boundary-matrix identities give single-block periodic-boundary constraints]
    \label{lem:block_diagonal_boundary_component_right_boundary_matrices}
    \lean{MPSTensor.blockDiagonal_boundary_component_chainGroundSpace_of_boundary_identities}
    \leanok
    \uses{lem:block_diagonal_boundary_crossing_right_boundary_matrix,
        lem:block_diagonal_boundary_component_crossing_windows_suffice}
    Assume $0<N$ and $L\le N$. Fix block-diagonal boundary conditions
    $X_j$. Suppose that, for every block $j$, every boundary-crossing
    interval beginning at $i$, and every configuration $\tau$, there is a
    matrix $Y_{j,i,\tau}$ such that, for every word $\beta$ of length
    $a=i+L-N$,
    \[
        \mu_j^N X_j A^j_\beta
        =
        A^j_\beta Y_{j,i,\tau} .
    \]
    Then,
    for every block $j$,
    \[
        \Gamma_N^{A_j}(\mu_j^NX_j)\in \mathcal G_{N,L}(A_j).
    \]
\end{lemma}

\begin{proof}
    \leanok
    \uses{lem:block_diagonal_boundary_crossing_right_boundary_matrix,
        lem:block_diagonal_boundary_component_crossing_windows_suffice}
    Lemma~\ref{lem:block_diagonal_boundary_crossing_right_boundary_matrix}
    turns the displayed boundary identity into the local constraint for each
    boundary-crossing interval.
    Lemma~\ref{lem:block_diagonal_boundary_component_crossing_windows_suffice}
    then yields
    \[
        \Gamma_N^{A_j}(\mu_j^NX_j)\in \mathcal G_{N,L}(A_j),
    \]
    handling both boundary-crossing intervals $N<i+L$ and non-crossing
    intervals $i+L\le N$.
\end{proof}

\begin{lemma}[Spanning matrix identities give single-block constraints]
    \label{lem:block_diagonal_boundary_component_complementary_word_identities}
    \lean{MPSTensor.blockDiagonal_boundary_component_chainGroundSpace_of_complementary_word_identities}
    \leanok
    \uses{thm:word_span_common_boundary_matrix,
        lem:block_diagonal_boundary_component_right_boundary_matrices}
    Assume $0<N$, $L\le N$, and that for every block $j$, the
    length-$N-L$ word products of $A_j$ span the full matrix algebra. Suppose
    that, for every block $j$, every boundary-crossing interval
    beginning at $i$, and every outside word $\rho$ of length $N-L$, there is
    a matrix $E_{j,i,\rho}$ such that for every word $\beta$ before the cut,
    of length $i+L-N$,
    \[
        \mu_j^N X_j A^j_\beta A^j_\rho
        =
        A^j_\beta E_{j,i,\rho}.
    \]
    Then, for every block $j$,
    \[
        \Gamma_N^{A_j}(\mu_j^NX_j)\in\mathcal G_{N,L}(A_j).
    \]
\end{lemma}

\begin{proof}
    \leanok
    \uses{thm:word_span_common_boundary_matrix,
        lem:block_diagonal_boundary_component_right_boundary_matrices}
    For fixed $j,i$, apply
    Theorem~\ref{thm:word_span_common_boundary_matrix} to the family
    \[
        Z_\beta=\mu_j^NX_jA^j_\beta,\qquad F_\beta=A^j_\beta.
    \]
    Since the length-$N-L$ outside-word products span the full matrix
    algebra, there is a single matrix $Y_{j,i}$ such that, for every $\beta$,
    \[
        \mu_j^NX_jA^j_\beta=A^j_\beta Y_{j,i} .
    \]
    Lemma~\ref{lem:block_diagonal_boundary_component_right_boundary_matrices}
    then gives the periodic-boundary constraints for each block.
\end{proof}

\begin{lemma}[Matrix identities under block injectivity give single-block constraints]
    \label{lem:block_diagonal_boundary_component_complementary_word_identities_injective}
    \lean{MPSTensor.blockDiagonal_boundary_component_chainGroundSpace_of_complementary_word_identities_of_injective}
    \leanok
    \uses{lem:block_diagonal_boundary_component_complementary_word_identities,
        thm:wordspan_propagation_unital}
    Assume $0<N$, $L\le N$, and $L+L_0\le N$. Suppose each block $A_j$
    is $L_0$-block-injective and satisfies the normalization equation
    \[
        \sum_a A^j_aA^{j\dagger}_a=I .
    \]
    Suppose also that the boundary-crossing matrix identities of
    Lemma~\ref{lem:block_diagonal_boundary_component_complementary_word_identities}
    hold for every outside word $\rho$ of length $N-L$. Then, for
    every block $j$,
    \[
        \Gamma_N^{A_j}(\mu_j^NX_j)\in\mathcal G_{N,L}(A_j).
    \]
\end{lemma}

\begin{proof}
    \leanok
    \uses{lem:block_diagonal_boundary_component_complementary_word_identities,
        thm:wordspan_propagation_unital}
    Since $L+L_0\le N$, the complementary length satisfies $L_0\le N-L$.
    Theorem~\ref{thm:wordspan_propagation_unital} therefore gives, for every
    block $j$,
    \[
        S_{N-L}(A_j)=\MN{D_j} .
    \]
    Apply
    Lemma~\ref{lem:block_diagonal_boundary_component_complementary_word_identities}.
\end{proof}

\begin{lemma}[Normalization for words]
    \label{lem:sum_evalWord_mul_conjTranspose_evalWord}
    \lean{MPSTensor.sum_evalWord_mul_conjTranspose_evalWord}
    \leanok
    \uses{}
    If
    \[
        \sum_a A_aA_a^\dagger=I ,
    \]
    then, for every length $M$,
    \[
        \sum_\rho A_\rho A_\rho^\dagger=I ,
    \]
    where $\rho$ ranges over all words of length $M$.
\end{lemma}

\begin{proof}
    \leanok
    \uses{}
    The case $M=0$ is the identity matrix.  For the induction step, write a
    word of length $M+1$ as $a\rho$.  Since $A_{a\rho}=A_aA_\rho$,
    \[
        \sum_{a,\rho} A_aA_\rho A_\rho^\dagger A_a^\dagger
        =
        \sum_a A_a
        \left(\sum_\rho A_\rho A_\rho^\dagger\right)
        A_a^\dagger
        =
        \sum_a A_aA_a^\dagger
        =
        I .
    \]
\end{proof}

\begin{lemma}[Compatibility hypotheses under block injectivity give single-block constraints]
    \label{lem:block_diagonal_boundary_component_pgvwc_comparison_injective}
    \lean{MPSTensor.blockDiagonal_boundary_component_chainGroundSpace_of_pgvwc_comparison_of_injective}
    \leanok
    \uses{lem:sum_evalWord_mul_conjTranspose_evalWord,
        lem:complementary_word_boundary_identities_from_pgvwc,
        lem:block_diagonal_boundary_component_complementary_word_identities_injective}
    Assume $0<N$, $L\le N$, and $L+L_0\le N$. Suppose each block $A_j$
    is $L_0$-block-injective and satisfies
    \[
        \sum_a A^j_aA^{j\dagger}_a=I .
    \]
    Fix block-diagonal boundary conditions $X_j$. For every boundary-crossing
    interval beginning at $i$, suppose there are matrices $C^j_{i,\rho}$,
    indexed by outside words $\rho$ of length $N-L$, such that for every word
    $\beta$ before the cut, of length $i+L-N$,
    \[
        A^j_\beta C^j_{i,\rho}
        =
        \bigl((\mu_j^NX_j)A^j_\beta\bigr)A^j_\rho .
    \]
    After opening the periodic boundary at the chosen cut, this is the source
    comparison
    \[
        A^j_{i_{m+1}}C^j_{i_1}=D^j_{i_{m+1}}A^j_{i_1},
        \qquad
        D^j_\beta=(\mu_j^NX_j)A^j_\beta .
    \]
    Then, for every block $j$,
    \[
        \Gamma_N^{A_j}(\mu_j^NX_j)\in\mathcal G_{N,L}(A_j).
    \]
\end{lemma}

\begin{proof}
    \leanok
    \uses{lem:sum_evalWord_mul_conjTranspose_evalWord,
        lem:complementary_word_boundary_identities_from_pgvwc,
        lem:block_diagonal_boundary_component_complementary_word_identities_injective}
    The normalization equation gives, by induction on word length,
    \[
        \sum_\rho A^j_\rho A^{j\dagger}_\rho=I
    \]
    for words $\rho$ of length $N-L$. Applying
    Lemma~\ref{lem:complementary_word_boundary_identities_from_pgvwc} with
    $X=\mu_j^NX_j$ gives, for every outside word $\rho$, a matrix
    $E_{j,i,\rho}$ such that for every word $\beta$ before the cut,
    \[
        \bigl((\mu_j^NX_j)A^j_\beta\bigr)A^j_\rho
        =
        A^j_\beta E_{j,i,\rho}.
    \]
    Lemma~\ref{lem:block_diagonal_boundary_component_complementary_word_identities_injective}
    then gives the periodic-boundary constraint for each block.
\end{proof}

\begin{lemma}[Boundary representation and crossing constraints give periodic components]
    \label{lem:block_diagonal_boundary_periodic_components_from_crossing_boundary}
    \lean{MPSTensor.exists_blockDiagonal_boundary_chainGroundSpace_of_crossing_pgvwc_comparison_of_boundary}
    \leanok
    \uses{thm:block_diagonal_boundary_crossing_pgvwc_comparison_chain_ground_space,
        lem:block_diagonal_boundary_component_pgvwc_comparison_injective}
    Assume $0<N$, $L\le N$, and $L+L_0\le N$. Suppose each block $A_j$
    is $L_0$-block-injective and satisfies
    \[
        \sum_a A^j_aA^{j\dagger}_a=I .
    \]
    Let
    \[
        \psi\in
        \mathcal G_{N,L}\!\left(\bigoplus_j\mu_jA_j\right),
        \qquad
        \psi=\Gamma_N^{\oplus_j\mu_jA_j}\!\left(\bigoplus_jX_j\right).
    \]
    If every boundary-crossing interval has the corresponding tail-word
    spanning property, then there are block-diagonal boundary conditions $X_j$
    with the displayed representation of $\psi$ and, for every $j$,
    \[
        \Gamma_N^{A_j}(\mu_j^NX_j)\in\mathcal G_{N,L}(A_j).
    \]
\end{lemma}

\begin{proof}
    \leanok
    \uses{thm:block_diagonal_boundary_crossing_pgvwc_comparison_chain_ground_space,
        lem:block_diagonal_boundary_component_pgvwc_comparison_injective}
    Theorem~\ref{thm:block_diagonal_boundary_crossing_pgvwc_comparison_chain_ground_space}
    supplies matrices $C^j_{i,\rho}$ satisfying
    \[
        A^j_\beta C^j_{i,\rho}
        =
        ((\mu_j^NX_j)A^j_\beta)A^j_\rho .
    \]
    Lemma~\ref{lem:block_diagonal_boundary_component_pgvwc_comparison_injective}
    then gives the single-block periodic constraints.
\end{proof}

\begin{lemma}[Boundary comparison implies trace comparison]
    \label{lem:block_diagonal_boundary_trace_decomposition_from_pgvwc_comparison}
    \lean{MPSTensor.blockDiagonal_boundary_trace_decomposition_of_pgvwc_comparison}
    \leanok
    Fix block-diagonal boundary matrices $X_j$. Suppose the opened-boundary
    matrices $C^j_{i,\rho}$ satisfy, for every boundary-crossing interval $i$,
    every outside word $\rho$, and every word $\beta$ before the cut,
    \[
        A^j_\beta C^j_{i,\rho}
        =
        ((\mu_j^NX_j)A^j_\beta)A^j_\rho .
    \]
    Then for every middle word $w$,
    \[
        \sum_j\tr\!\left(A^j_\beta C^j_{i,\rho}A^j_w\right)
        =
        \sum_j
        \tr\!\left(((\mu_j^NX_j)A^j_\beta)A^j_\rho A^j_w\right).
    \]
\end{lemma}

\begin{proof}
    \leanok
    Substituting the displayed comparison identity in each summand gives
    \[
        \tr\!\left(A^j_\beta C^j_{i,\rho}A^j_w\right)
        =
        \tr\!\left(((\mu_j^NX_j)A^j_\beta)A^j_\rho A^j_w\right).
    \]
    Summing over \(j\) gives the conclusion.
\end{proof}

\begin{lemma}[Boundary trace comparisons give single-block constraints]
    \label{lem:block_diagonal_boundary_component_trace_decomposition_injective}
    \lean{MPSTensor.blockDiagonal_boundary_component_chainGroundSpace_of_trace_decomposition_of_injective}
    \leanok
    \uses{lem:block_boundary_cde_identities_from_trace_decompositions,
        lem:block_diagonal_boundary_component_complementary_word_identities_injective}
    Assume $0<N$, $L\le N$, and $L+L_0\le N$. Suppose each block $A_j$
    is $L_0$-block-injective and satisfies
    \[
        \sum_a A^j_aA^{j\dagger}_a=I .
    \]
    Assume also that the simultaneous tuples $w\mapsto(A^1_w,\ldots,A^g_w)$
    at the middle-word length $m$ span the full product algebra
    $\prod_j\MN{D_j}$.
    Fix complex numbers $\mu_j$ and block-diagonal boundary conditions
    $X_j$. For every boundary-crossing interval beginning at $i$, suppose
    there are matrices $C^j_{i,\rho}$, indexed by outside words
    $\rho$ of length $N-L$, such that for every word $\beta$ before the cut,
    outside word $\rho$, and middle word $w$,
    \[
        \sum_j\tr\!\left(A^j_\beta C^j_{i,\rho}A^j_w\right)
        =
        \sum_j
        \tr\!\left(((\mu_j^NX_j)A^j_\beta)A^j_\rho A^j_w\right).
    \]
    Then, for every block $j$,
    \[
        \Gamma_N^{A_j}(\mu_j^NX_j)\in\mathcal G_{N,L}(A_j).
    \]
\end{lemma}

\begin{proof}
    \leanok
    \uses{lem:block_boundary_cde_identities_from_trace_decompositions,
        lem:block_diagonal_boundary_component_complementary_word_identities_injective}
    Applying
    Lemma~\ref{lem:block_boundary_cde_identities_from_trace_decompositions}
    with the same scalars $\mu_j$ and boundary matrices $X_j$, for every
    block $j$ and every boundary-crossing interval $i$, this gives a matrix
    $E^j_i$ such that, for every word $\beta$ before the cut,
    \[
        (\mu_j^NX_j)A^j_\beta=A^j_\beta E^j_i .
    \]
    Hence for every outside word $\rho$, taking
    $E_{j,i,\rho}=E^j_iA^j_\rho$ gives
    \[
        ((\mu_j^NX_j)A^j_\beta)A^j_\rho=A^j_\beta E_{j,i,\rho} .
    \]
    These are precisely the block-diagonal boundary-condition equations of
    Lemma~\ref{lem:block_diagonal_boundary_component_complementary_word_identities_injective},
    which gives the periodic-boundary constraint for each block.
\end{proof}

\begin{lemma}[Boundary trace comparisons with a block-diagonal representation]
    \label{lem:block_diagonal_boundary_periodic_components_from_trace_decomposition_boundary}
    \lean{MPSTensor.exists_blockDiagonal_boundary_chainGroundSpace_of_trace_decomposition_of_boundary}
    \leanok
    \uses{lem:block_diagonal_boundary_component_trace_decomposition_injective}
    Assume $0<N$, $L\le N$, and $L+L_0\le N$. Suppose each block $A_j$
    is $L_0$-block-injective and satisfies
    \[
        \sum_a A^j_aA^{j\dagger}_a=I .
    \]
    Assume also that the simultaneous tuples $w\mapsto(A^1_w,\ldots,A^g_w)$
    at the middle-word length $m$ span the full product algebra
    $\prod_j\MN{D_j}$. Let
    \[
        \psi=\Gamma_N^{\oplus_j\mu_jA_j}\!\left(\bigoplus_jX_j\right)
    \]
    for some block-diagonal boundary conditions $X_j$. Suppose that whenever
    $\psi$ has such a representation, there are matrices $C^j_{i,\rho}$
    such that, for every boundary-crossing interval $i$, every outside word
    $\rho$ of length $N-L$, every word $\beta$ of length $i+L-N$ before the
    cut, and every word $w$ of length $m$,
    \[
        \sum_j\tr\!\left(A^j_\beta C^j_{i,\rho}A^j_w\right)
        =
        \sum_j
        \tr\!\left(((\mu_j^NX_j)A^j_\beta)A^j_\rho A^j_w\right).
    \]
    Then there are block-diagonal boundary conditions $X_j$ with
    \[
        \psi=\Gamma_N^{\oplus_j\mu_jA_j}\!\left(\bigoplus_jX_j\right)
    \]
    and, for every $j$,
    \[
        \Gamma_N^{A_j}(\mu_j^NX_j)\in\mathcal G_{N,L}(A_j).
    \]
\end{lemma}

\begin{proof}
    \leanok
    \uses{lem:block_diagonal_boundary_component_trace_decomposition_injective}
    Choose a block-diagonal boundary representation of $\psi$. For that
    representation, the boundary trace comparison supplies the matrices
    $C^j_{i,\rho}$. Applying
    Lemma~\ref{lem:block_diagonal_boundary_component_trace_decomposition_injective}
    gives the desired containment
    $\Gamma_N^{A_j}(\mu_j^NX_j)\in\mathcal G_{N,L}(A_j)$ for every block $j$.
\end{proof}
